
To quantify the clustering of alignments within an Ising magnet, we define a quantity called the \textbf{correlation function}, $c(r)$. Take any two dipoles, $i$ and $j$, separated by a distance $r$, and compute the product of their states: $s_i s_j$. This product is 1 if the dipoles are parallel and -1 if the dipoles are antiparallel. Now average this quantity over \textit{all} pairs that are separated by a fixed distance $r$, to obtain a measure of the tendency of dipoles to be "correlated" over this distance. Finally, to remove the effect of any overall magnetization of the system, subtract off the square of the average $s$. Written as an equation, then, the correlation function is

\begin{equation*}
	c(r) = \overline{s_i s_j} - \overline{s_i}^2
\end{equation*}

where it is understood that the first term averages over all pairs at the fixed distance $r$. Technically, the averages should also be taken over all possible states of the system, but don't do this yet. 

\textbf{(a)} Add a routine to the \texttt{ising} program to compute the correlation function for the current state of the lattice, averaging over all pairs separated either vertically or horizontally (but not diagonally) by $r$ units of distance, where $r$ varies from 1 to half the lattice size. Have the program execute this routine periodically and plot the results as a bar graph.

\textbf{Answer:} I ran the correlation for a lattice of size 100 and $T=2.27$ and created the following bar graph:

\begin{figure}[H]
	\centering
	\includegraphics[scale=0.75]{corrA.png}
\end{figure}

We can see the function decays exponentially as $r$ increases and after around $r=20$ it is roughly zero.

\textbf{(b)} Run this program at a variety of temperatures above, below, and near the critical point. Use a lattice size of at least 20, preferably larger (especially near the critical point). Describe the behavior of the correlation function at each temperature.

\textbf{Answer:} I created charts of the correlation for the following temperatures: 0.5, 1, 2, 2.27 (the critical temperature), 3, and 5, using a lattice size of 50 so that equilibrium can be reached quicker:

\begin{figure}[H]
	\centering
	
	\begin{subfigure}{0.32\linewidth}
		\includegraphics[width=\linewidth]{corrB_0.5.png}
	\end{subfigure}
	\hfill
	\begin{subfigure}{0.32\linewidth}
		\includegraphics[width=\linewidth]{corrB_1.png}
	\end{subfigure}
	\hfill
	\begin{subfigure}{0.32\linewidth}
		\includegraphics[width=\linewidth]{corrB_2.png}
	\end{subfigure}
	
	\vspace{0.5em}
	
	\begin{subfigure}{0.32\linewidth}
		\includegraphics[width=\linewidth]{corrB_Tc.png}
	\end{subfigure}
	\hfill
	\begin{subfigure}{0.32\linewidth}
		\includegraphics[width=\linewidth]{corrB_3.png}
	\end{subfigure}
	\hfill
	\begin{subfigure}{0.32\linewidth}
		\includegraphics[width=\linewidth]{corrB_5.png}
	\end{subfigure}
\end{figure}

We can see that t temperatures below $T_c$, the correlation function is small in absolute terms, even at $r=1$, but correlation continues to large distances. This makes sense because the lattice is rather homogeneous at equilibrium. Above $T_c$, the correlation function drops to 0 rapidly because there are at most small structures, and the magnetization as a whole is largely uncorrelated.

\textbf{(c)} Now add code to compute the \textit{average} correlation function over the duration of a run. (However, it's best to let the system "equilibrate" to a typical state before you begin accumulating averages.) The \textbf{correlation length} as defined as the distance over which the correlation function decreases by a factor of $e$. Estimate the correlation length at each temperature, and plot a graph of the correlation length vs. \texttt{T}.

\textbf{Answer:} I opted to sample correlation functions from 50 of the last 1000 frames of the simulation and take the average of them. By converting to a scatter plot and log scale, I can come up with a linear fit to estimate the correlation length, which should be the negative reciprocal of the slope:

\begin{figure}[H]
	\centering
	\includegraphics[scale=0.75]{corrC_1.png}
\end{figure}

I was able to plot correlation lengths for a series of temperatures surrounding the critical temperature, with a lattice size of 30:

\begin{figure}[H]
	\centering
	\includegraphics[scale=0.75]{corrC_2.png}
\end{figure}

While this data is a bit noisy, it can be seen that the correlation length sharply increases around $T_c$. Above this temperature, large scale structures are broken up so dipoles are uncorrelated. Below this temperature, magnetization is unified at equilibrium, so there is no meaningful correlation to be ascertained. Only around the critical temperature do these effects balance out and yield a significant correlation length, signifying the unique structures found at $T_c$.